\subsection{\texorpdfstring{\(\mathfrak{sl}_2\) -TRIPLETS IN SEMI-SIMPLE LIE ALGEBRAS}{\textbackslash mathfrak\{sl\}\_2 -TRIPLETS IN SEMI-SIMPLE LIE ALGEBRAS}}

Lemma 5. Let V be a finite dimensional vector space, A and B endomorphisms of V. Assume that A is nilpotent and that \([A, [A, B]] = 0\) . Then AB is nilpotent.

Put C = {[}A, B{]}. Since {[}A, C{]} = 0,

\[
[ \mathrm{A}, \mathrm{BC} ^ {p} ] = [ \mathrm{A}, \mathrm{B} ] \mathrm{C} ^ {p} = \mathrm{C} ^ {p + 1}
\]

for every integer \(p \geq 0\) . Consequently, \(\mathrm{Tr}(\mathrm{C}^{p}) = 0\) for \(p \geq 1\) , which proves that C is nilpotent (Algebra, Chap. VII, §3, no. 5, Cor. 4 of Prop. 13). Now let \(\bar{k}\) be an algebraic closure of k, and let \(\lambda \in \bar{k}\) , \(x \in V \otimes_{k} \bar{k}\) be such that ABx = \(\lambda x\) , \(x \neq 0\) . The relation \([[B, A], A] = 0\) shows that \([B, A^{p}] = p[B, A]A^{p-1}\) for every integer \(p \geq 0\) . Let r be the smallest integer such that \(A^{r}x = 0\) . Then

\[
\lambda \mathrm{A} ^ {r - 1} x = \mathrm{A} ^ {r - 1} \mathrm{AB} x = \mathrm{A} ^ {r} \mathrm{B} x = \mathrm{BA} ^ {r} x - [ \mathrm{B}, \mathrm{A} ^ {r} ] x = - r [ \mathrm{B}, \mathrm{A} ] \mathrm{A} ^ {r - 1} x.
\]

Since {[}B, A{]} is nilpotent and since \(A^{r-1}x \neq 0\) , this proves that \(\lambda = 0\) . Thus, all the eigenvalues of AB are zero, hence the lemma.

Lemma 6. Let \(h, x \in \mathfrak{g}\) be such that \([h, x] = 2x\) and \(h \in (\operatorname{ad} x)(\mathfrak{g})\) . Then there exists \(y \in \mathfrak{g}\) such that \((x, h, y)\) is either \((0, 0, 0)\) or an \(\mathfrak{sl}_2\) -triplet.

Let \(\mathfrak{g}'\) be the solvable Lie algebra \(kh + kx\) . Since \(x \in [\mathfrak{g}', \mathfrak{g}']\) , \(\operatorname{ad}_{\mathfrak{g}} x\) is nilpotent (Chap. I, §5, no. 3, Th. 1); let \(\mathfrak{n}\) be its kernel. Since {[}ad \(h\) , ad \(x] = 2\) ad \(x\) , we have (ad \(h\) ) \(\mathfrak{n} \subset \mathfrak{n}\) . Let \(z \in \mathfrak{g}\) be such that \(h = -[x, z]\) . For any integer \(n \geq 0\) , put \(\mathrm{M}_n = (\mathrm{ad}~x)^n\mathfrak{g}\) . If \(n > 0\) , we have (§1, no. 1, Lemma 1)

\[
[ \mathrm{ad} z, (\mathrm{ad} x) ^ {n} ] = n ((\mathrm{ad} h) - n + 1) (\mathrm{ad} x) ^ {n - 1}
\]

so, if \(u \in M_{n-1}\) ,

\[
n ((\operatorname{ad} h) - n + 1) u \in (\operatorname{ad} z) (\operatorname{ad} x) u + \mathrm{M} _ {n}.
\]

Since (ad \(h)\mathfrak{n}\subset \mathfrak{n}\) , it follows that

\[
((\operatorname{ad} h) - n + 1) (\mathfrak {n} \cap \mathrm{M} _ {n - 1}) \subset \mathfrak {n} \cap \mathrm{M} _ {n}.
\]

Since ad x is nilpotent, \(M_{n} = 0\) for sufficiently large n.~Consequently, the eigenvalues of ad \(h|n\) are integers \(\geq 0\) . Thus, the restriction of ad \(h + 2\) to n is invertible.

Now \([h,z] + 2z\in \mathfrak{n}\) since

\[
\begin{array}{r} [ x, [ h, z ] + 2 z ] = [ [ x, h ], z ] + [ h, [ x, z ] ] + 2 [ x, z ] \\ = [ - 2 x, z ] + [ h, - h ] + 2 [ x, z ] = 0. \end{array}
\]

Hence there exists \(z' \in \mathfrak{n}\) such that \([h, z'] + 2z' = [h, z] + 2z\) , that is, \([h, y] = -2y\) , putting \(y = z - z'\) . Since \([x, y] = [x, z] = -h\) , this completes the proof.

PROPOSITION 2 (Jacobson-Morozov). Assume that \(\mathfrak{g}\) is semi-simple. Let \(x\) be a non-zero nilpotent element of \(\mathfrak{g}\) . There exist \(h, y \in \mathfrak{g}\) such that \((x, h, y)\) is an \(\mathfrak{sl}_2\) -triplet.

Let \(\mathfrak{n}=\operatorname{Ker}(\operatorname{ad}x)^{2}\) . If \(z\in\mathfrak{n}\) , then {[}ad x, {[}ad x, ad z{]}{]} = ad({[}x, {[}x, z{]}{]}) = 0. By Lemma 5, ad \(x\circ ad\) z is nilpotent, so \(\operatorname{Tr}(\operatorname{ad}x\circ\operatorname{ad}z)=0\) . This shows that x is orthogonal to n with respect to the Killing form \(\Phi\) of g. Since

\[
\Phi ((\operatorname{ad} x) ^ {2} y, y ^ {\prime}) = \Phi (y, (\operatorname{ad} x) ^ {2} y ^ {\prime})
\]

for all \(y, y' \in g\) , and since \(\Phi\) is non-degenerate, the orthogonal complement of n is the image of \((\operatorname{ad} x)^{2}\) . Hence there exists \(y' \in g\) such that \(x = (\operatorname{ad} x)^{2}y'\) . Put

\[
h = - 2 [ x, y ^ {\prime} ];
\]

we have \([h,x]=2x\) and \(h\in(\operatorname{ad}x)(\mathfrak{g})\) . It now suffices to apply Lemma 6.

COROLLARY. Assume that \(\mathfrak{g}\) is semi-simple. Let \(\mathrm{G}\) be a group of automorphisms of \(\mathfrak{g}\) containing \(\operatorname{Aut}_e(\mathfrak{g})\) . The map which associates to any \(\mathfrak{sl}_2\) -triplet \((x,h,y)\) in g the nilpotent element x defines, by passage to the quotient, a bijection from the set of G-conjugacy classes of \(sl_{2}\) -triplets to the set of G-conjugacy classes of non-zero nilpotent elements.

This follows from Prop. 1 and 2.

Lemma 7. Let K be a commutative field with at least 4 elements. Let G be the group of matrices \(\begin{pmatrix}\alpha&\beta\\0&\alpha^{-1}\end{pmatrix}\) where \(\alpha\in K^{*},\beta\in K\) . Let \(G'\) be the group of such matrices such that \(\alpha=1\) . Then \(G'=(G,G)\) .

If \(\alpha, \alpha' \in \mathrm{K}^*\) and \(\beta, \beta' \in \mathrm{K}\) ,

\[
\begin{array}{c} \left( \begin{array}{c c} \alpha & \beta \\ 0 & \alpha^ {- 1} \end{array} \right) \left( \begin{array}{c c} \alpha^ {\prime} & \beta^ {\prime} \\ 0 & \alpha^ {- 1} \end{array} \right) \left( \begin{array}{c c} \alpha & \beta \\ 0 & \alpha^ {- 1} \end{array} \right) ^ {- 1} \left( \begin{array}{c c} \alpha^ {\prime} & \beta^ {\prime} \\ 0 & \alpha^ {- 1} \end{array} \right) ^ {- 1} \\ = \left( \begin{array}{c c} 1 & - \alpha^ {\prime} \beta^ {\prime} - \alpha \beta \alpha^ {2} + \alpha^ {2} \alpha^ {\prime} \beta^ {\prime} + \alpha \beta \\ 0 & 1 \end{array} \right). \end{array}
\]

In particular,

\[
\left( \begin{array}{c c} 1 & \beta \\ 0 & 1 \end{array} \right) \left( \begin{array}{c c} \alpha^ {\prime} & 0 \\ 0 & \alpha^ {- 1} \end{array} \right) \left( \begin{array}{c c} 1 & \beta \\ 0 & 1 \end{array} \right) ^ {- 1} \left( \begin{array}{c c} \alpha^ {\prime} & 0 \\ 0 & \alpha^ {- 1} \end{array} \right) ^ {- 1} = \left( \begin{array}{c c} 1 & \beta (1 - \alpha^ {2}) \\ 0 & 1 \end{array} \right).
\]

But there exists \(\alpha_0' \in \mathrm{K}^*\) such that \(\alpha_0' \neq 1\) and \(\alpha_0' \neq -1\) , and then \(k.(1 - \alpha_0'^2) = k\) , hence the lemma.

PROPOSITION 3. Assume that \(\mathfrak{g}\) is semi-simple. The group \(\mathrm{Aut}_e(\mathfrak{g})\) is equal to its derived group. If \(\mathfrak{g}\) is splittable, \(\mathrm{Aut}_e(\mathfrak{g})\) is the derived group of \(\mathrm{Aut}_0(\mathfrak{g})\) .

Let \(x\) be a non-zero nilpotent element of \(\mathfrak{g}\) . Choose \(h, y \in \mathfrak{g}\) be such that \((x, h, y)\) is an \(\mathfrak{sl}_2\) -triplet (Prop. 2). The subalgebra \(\mathfrak{s}\) of \(\mathfrak{g}\) generated by \((x, h, y)\) can be identified with \(\mathfrak{sl}(2, k)\) . Let \(\rho\) be the representation \(z \mapsto \operatorname{ad}_{\mathfrak{g}} z\) of \(\mathfrak{s} = \mathfrak{sl}(2, k)\) on \(\mathfrak{g}\) , and let \(\pi\) be the representation of \(\mathbf{SL}(2, k)\) compatible with \(\rho\) ( \(\S 1\) , no. 4). The image of \(\pi\) is generated by the \(\exp(t\operatorname{ad}_{\mathfrak{g}} x)\) and the \(\exp(t\operatorname{ad}_{\mathfrak{g}} y)\) with \(t \in k\) (Algebra, Chap. III, \(\S 8\) , no. 9, Prop. 17), hence is contained in \(\operatorname{Aut}_e(\mathfrak{g})\) . Since \(\mathbf{SL}(2, k)\) is equal to its derived group (Lemma 7 and loc. cit.), \(\exp(\operatorname{ad}_{\mathfrak{g}} x)\) belongs to the derived group G of \(\operatorname{Aut}_e(\mathfrak{g})\) . Hence \(\operatorname{Aut}_e(\mathfrak{g})\) is equal to G. Assume now that \(\mathfrak{g}\) is splittable. Since \(\operatorname{Aut}_0(\mathfrak{g}) / \operatorname{Aut}_e(\mathfrak{g})\) is commutative ( \(\S 5\) , no. 3, Remark 3), the preceding proves that the derived group of \(\operatorname{Aut}_0(\mathfrak{g})\) is \(\operatorname{Aut}_e(\mathfrak{g})\) .

